\chapter{数量级：测量宇宙}\label{ch:g10:orders-of-magnitude}

物理学始于观点止步之处：始于测量。但物理所测量的对象
跨越的范围大得惊人——原子核的直径约为
\qty{e-15}{m}，可观测宇宙的直径约为 \qty{e27}{m}，两者相差
$10^{42}$ 倍，远超任何心理图像所能容纳的范围。
本章介绍驯服这一范围的工具——\omterm{def:g10:orders-of-magnitude:unit}{单位}、
十的幂、\omterm{def:g10:orders-of-magnitude:sigfig}{有效数字}——以及此后每一页都会用到的两个习惯：
写出\omterm{def:g10:orders-of-magnitude:unit}{单位}，检查\omterm{def:g10:orders-of-magnitude:oom}{数量级}。

\section{单位：测量的语法}

\begin{definition}[物理量与国际单位制]\label{def:g10:orders-of-magnitude:unit}
称\emph{物理量}\index{物理量}为
世界上可以被测量的某种属性，例如长度、时长、质量。
测量一个物理量，就是数一数某个参考量——
\emph{单位}\index{单位}——在其中包含几次。
\emph{国际单位制}\index{国际单位制 (SI)}
（SI）规定了七个基本单位，其中三个贯穿全书：
长度用\emph{米}\index{米}（\unit{m}），
质量用\emph{千克}\index{千克}（\unit{kg}），
时间用\emph{秒}\index{秒}（\unit{s}）。（其余四个
——安培、开尔文、摩尔、坎德拉——将在电学、热学
和化学内容出现时再引入。）
\end{definition}

\begin{remark}[没有单位的数不是结果]\label{rem:g10:orders-of-magnitude:units}
“距离是 $384$”毫无意义：\qty{384}{mm}，
\qty{384}{km}、\qty{384}{Mm}是三种完全不同的距离（最后一个
恰好是地月距离）。1999 年，美国国家航空航天局（NASA）损失了
价值 \$327 百万的\emph{火星气候轨道器}，原因在于一个程序以
磅力计算推进器推力，而另一个程序却按牛顿读取。
本书约定，自此行起直至全书结束，凡测得或算得的数值
在计算的每一行中都须带上\omterm{def:g10:orders-of-magnitude:unit}{单位}。
\end{remark}

\begin{definition}[国际单位制词头]\label{def:g10:orders-of-magnitude:prefixes}
称附在\omterm{def:g10:orders-of-magnitude:unit}{单位}前的\emph{词头}\index{词头 (SI)}，
将该\omterm{def:g10:orders-of-magnitude:unit}{单位}乘以一个固定的十的幂次：
\begin{center}
\begin{tabular}{llc@{\qquad\qquad}llc}
飞 & f & $10^{-15}$ & 厘 & c & $10^{-2}$ \\
皮 & p & $10^{-12}$ & 千 & k & $10^{3}$  \\
纳 & n & $10^{-9}$  & 兆 & M & $10^{6}$  \\
微 & $\mu$ & $10^{-6}$ & 吉 & G & $10^{9}$ \\
毫 & m & $10^{-3}$  & 太 & T & $10^{12}$ \\
\end{tabular}
\end{center}
由此可得 $\qty{1}{nm} = \qty{e-9}{m}$ 以及 $\qty{1}{km} = \qty{e3}{m}$。
\end{definition}

\begin{example}[读懂词头]\label{ex:g10:orders-of-magnitude:prefixes}
绿光的波长约为 $\qty{500}{nm} = \qty{5.00e-7}{m}$；人的头发
粗细约为 $\qty{70}{\micro m} = \qty{7.0e-5}{m}$；月球绕地球
运行的距离约为 $\qty{384}{Mm} = \qty{3.84e8}{m}$。一滴雨滴
的质量约为 $\qty{30}{mg} = \qty{3.0e-5}{kg}$——请注意，\omterm{def:g10:orders-of-magnitude:unit}{国际单位制}的
质量基本\omterm{def:g10:orders-of-magnitude:unit}{单位}是七个基本\omterm{def:g10:orders-of-magnitude:unit}{单位}中唯一自带\omterm{def:g10:orders-of-magnitude:prefixes}{词头}的那一个。
\end{example}

\section{科学记数法与有效数字}

\begin{definition}[科学记数法]\label{def:g10:orders-of-magnitude:scinot}
一个数若写成\emph{科学记数法}\index{科学记数法}形式，
则可写作：
\[
a \times 10^{n}, \qquad 1 \leq a < 10, \quad n \text{ 为整数.}
\]
因子 $a$ 称为\emph{尾数}\index{尾数}，整数
$n$ 称为\emph{指数}\index{指数}。
由此可得 $\num{149600000} = \num{1.496e8}$ 与
$\num{0.00052} = \num{5.2e-4}$。
\end{definition}

\begin{proposition}[十的幂的运算]\label{prop:g10:orders-of-magnitude:powers}
对所有整数 $m$ 与 $n$：
\[
10^{m} \times 10^{n} = 10^{m+n}, \qquad
\frac{1}{10^{n}} = 10^{-n}, \qquad
\left(10^{m}\right)^{n} = 10^{mn}.
\]
\end{proposition}

\begin{proof}[部分证明]
对正\omterm{def:g10:orders-of-magnitude:scinot}{指数}，数一数因子个数：$10^{m} \times 10^{n}$ 是一个
由 $m + n$ 个十相乘的积，而 $(10^m)^n$ 则是 $n$ 组、每组 $m$ 个十相乘。
规定 $10^{-n}$ 即为 $1/10^{n}$，这正是为了
使第一条规则对所有整数都成立：它要求
$10^{n} \times 10^{-n} = 10^{0} = 1$。配套的数学分册
将全面研究\omterm{def:g10:orders-of-magnitude:scinot}{指数}法则。
\end{proof}

\begin{example}[用科学记数法计算]\label{ex:g10:orders-of-magnitude:scinot}
按质量计，太阳内能容纳多少个地球？取
$m_{\text{太阳}} = \qty{1.99e30}{kg}$ 与
$m_{\text{地球}} = \qty{5.97e24}{kg}$：
\[
\frac{\num{1.99e30}}{\num{5.97e24}}
= \frac{1.99}{5.97} \times 10^{30-24}
= 0.333 \times 10^{6}
= \num{3.33e5}.
\]
\omterm{def:g10:orders-of-magnitude:scinot}{尾数}与\omterm{def:g10:orders-of-magnitude:scinot}{指数}分别处理；最终结果须重新规范化，
使\omterm{def:g10:orders-of-magnitude:scinot}{尾数}重新落在 $[1, 10)$ 内。
\end{example}

\begin{definition}[有效数字]\label{def:g10:orders-of-magnitude:sigfig}
称测得数的\emph{有效数字}\index{有效数字}为
测量所真正可靠的数字，即除前导零以外
的全部写出的数字。末尾的零也计入——
$\qty{5.20e-3}{s}$（三位有效数字）
所表示的精度高于 $\qty{5.2e-3}{s}$（两位）。\omterm{def:g10:orders-of-magnitude:scinot}{科学记数法}
一眼即可看出位数：\omterm{def:g10:orders-of-magnitude:scinot}{尾数}恰好包含全部
有效数字。
\end{definition}

\begin{method}[保留几位数字]\label{met:g10:orders-of-magnitude:sigfig}
计算结果的精度不能高于其原始数据的精度。
\begin{enumerate}
  \item 在乘积或商中，保留的\omterm{def:g10:orders-of-magnitude:sigfig}{有效数字}位数应
        与\emph{精度最低}的因子保持相同。
  \item 只对最终结果取整，中间步骤不取整。
  \item 拿不准时，三位\omterm{def:g10:orders-of-magnitude:sigfig}{有效数字}即为本书的
        工作精度。
\end{enumerate}
\end{method}

\begin{example}[地球的周长]\label{ex:g10:orders-of-magnitude:sigfig}
地球的直径为 $d = \qty{1.27e7}{m}$（三位\omterm{def:g10:orders-of-magnitude:sigfig}{有效数字}），
因此其周长为
$\pi d = \num{3.14159} \times \qty{1.27e7}{m} = \qty{3.99e7}{m}$ ---
三位\omterm{def:g10:orders-of-magnitude:sigfig}{有效数字}，而非计算器显示的八位。写成
$\qty{39898227}{m}$ 则等于声称能将这颗行星精确到米。
\end{example}

\begin{omfigure}
\begin{tikzpicture}[x=1.9cm, font=\small]
  \draw[-Stealth] (-0.3,0) -- (5.45,0);
  \foreach \e in {0,...,5} {
    \draw (\e,-0.1) -- (\e,0.1);
    \node[below] at (\e,-0.12) {$10^{\e}$};
  }
  \foreach \x in {2.301,2.477,2.602,2.699,2.778,2.845,2.903,2.954}
    \draw (\x,-0.05) -- (\x,0.05);
  \draw[Stealth-Stealth, omMeth] (0,0.45) -- (1,0.45)
    node[midway, above] {$\times 10$};
  \fill[omDef] (2.477,0) circle (1.7pt);
  \node[above=3pt, omDef] at (2.477,0.06) {$3 \times 10^{2}$};
  \draw[-Stealth, omDef, dashed] (2.38,0.22) -- (2.06,0.22);
  \fill[omThm] (3.845,0) circle (1.7pt);
  \node[above=3pt, omThm] at (3.845,0.06) {$7 \times 10^{3}$};
  \draw[-Stealth, omThm, dashed] (3.90,0.22) -- (3.96,0.22);
\end{tikzpicture}

{\small 刻度尺上每一等距的刻度都\emph{乘以}十。小刻线为
$200, 300, \dots, 900$：每段内挤满了整整一个十倍的范围。
在此刻度尺上 $\num{300}$ 靠近 $10^{2}$，而
$\num{7000}$ 靠近 $10^{4}$——即它们的\omterm{def:g10:orders-of-magnitude:oom}{数量级}。}
\end{omfigure}

\section{宇宙的阶梯}

\begin{definition}[数量级]\label{def:g10:orders-of-magnitude:oom}
称某一量的\emph{数量级}\index{数量级}
为与其最接近的十的幂：将该量写成
 $a \times 10^{n}$，其中 $1 \leq a < 10$，取 $10^{n}$，当 $a < 5$ 时，
$10^{n+1}$，当 $a \geq 5$ 时。（分界处的确切约定
从不影响结论：数量级的表述只在
十倍之内有意义。）
\end{definition}

\begin{example}[数量级举例]\label{ex:g10:orders-of-magnitude:oom}
人的一生约持续 $80$ 年，即\
$80 \times \qty{3.16e7}{s} \approx \qty{2.5e9}{s}$：\omterm{def:g10:orders-of-magnitude:oom}{数量级}为
$\qty{e9}{s}$——约数十亿秒，要好好利用。人的
身高为 $\qty{1.7}{m}$，\omterm{def:g10:orders-of-magnitude:oom}{数量级}为 $\qty{e0}{m}$。珠穆朗玛峰，
$\qty{8848}{m}$：\omterm{def:g10:orders-of-magnitude:oom}{数量级}为 $\qty{e4}{m}$，因为 $8.848 \geq 5$。
\end{example}

将这些估算按\omterm{def:g10:orders-of-magnitude:scinot}{指数}排列，便得到本章的总图
景：一座阶梯，每一级都是
比下一级宽十倍的世界。

\begin{omfigure}
\begin{tikzpicture}[x=0.26cm, font=\small]
  \draw[-Stealth] (-16.5,0) -- (28.7,0);
  \foreach \e in {-15,-10,-5,0,5,10,15,20,25} {
    \draw (\e,-0.12) -- (\e,0.12);
    \node[below, font=\footnotesize] at (\e,-0.14) {$10^{\e}$};
  }
  \foreach \e/\t in {-15/原子核, -10/原子, -6/细菌, 0/人,
                     7/地球, 11/日地距离, 16/光年,
                     21/银河系, 27/宇宙} {
    \draw[omDef, thick] (\e,0.14) -- (\e,0.62);
    \node[rotate=52, anchor=south west, inner sep=1pt, omDef]
      at (\e,0.66) {\t};
  }
  \node[below, font=\footnotesize] at (26.6,-0.75) {尺寸（\unit{m}）};
\end{tikzpicture}

{\small 宇宙的阶梯，以米为\omterm{def:g10:orders-of-magnitude:unit}{单位}：$42$ 个十的幂
将原子核与可观测宇宙隔开，而人类
位于距原子十级、距顶端二十七级之处。}
\end{omfigure}

在地球这一级之上，米已不便于使用，天文学
保留了自己的\omterm{def:g10:orders-of-magnitude:unit}{单位}，两者皆由这座阶梯本身定义。

\begin{definition}[天文单位与光年]\label{def:g10:orders-of-magnitude:lightyear}
称\emph{天文单位}\index{天文单位}
（$\qty{1}{au} = \qty{1.496e11}{m}$）为地球与太阳之间的平均
距离。称\emph{光年}\index{光年}（\unit{ly}）为
光在一年中所传播的距离。光在真空中以
\[
c = \qty{3.00e8}{m/s},
\]
即自然界允许的最快速率——因此光年是一个距离，
而不是时间。
\end{definition}

\begin{omfigure}
\begin{tikzpicture}[x=1.08cm, font=\small]
  \draw[-Stealth] (7.6,0) -- (17.5,0);
  \foreach \e in {8,10,12,14,16} {
    \draw (\e,-0.12) -- (\e,0.12);
    \node[below, font=\footnotesize] at (\e,-0.14) {$10^{\e}$};
  }
  \foreach \x/\t/\d/\a in {8.58/月球/{\qty{1.3}{s}}/north,
                           11.17/太阳/{\qty{8.3}{min}}/north,
                           12.65/海王星/{\qty{4.2}{h}}/north,
                           15.98/{$\qty{1}{ly}$}/{\qty{1}{yr}}/north east,
                           16.60/比邻星/{\qty{4.2}{yr}}/north west} {
    \draw[omDef, thick] (\x,0.14) -- (\x,0.55);
    \node[rotate=52, anchor=south west, inner sep=1pt, omDef]
      at (\x,0.58) {\t};
    \node[anchor=\a, font=\footnotesize, omThm] at (\x,-0.55) {\d};
  }
  \node[right, font=\footnotesize] at (17.5,0) {距离（\unit{m}）};
\end{tikzpicture}

{\small 乘法标尺上标出了各天体距地球的距离，每个距离下方以红色标出光的
传播时间：月球以光秒计，
行星以光时计，恒星以\omterm{def:g10:orders-of-magnitude:lightyear}{光年}计。}
\end{omfigure}

\begin{example}[光年的米数]\label{ex:g10:orders-of-magnitude:ly}
一年包含 $365.25 \times 24 \times 3600 = \qty{3.156e7}{s}$
（恰好接近 $\pi \times 10^{7}\,\unit{s}$）。由此
\[
\qty{1}{ly} = c \times \qty{3.156e7}{s}
= \qty{3.00e8}{m/s} \times \qty{3.156e7}{s}
= \qty{9.47e15}{m}.
\]
比邻星是除太阳外离我们最近的恒星，位于
$\qty{4.2}{ly} \approx \qty{4.0e16}{m}$ 远处。光穿越一个
日地距离只需八分钟，而光从那里
到达地球则需四年以上。
\end{example}

\begin{method}[数量级的合理性检验]\label{met:g10:orders-of-magnitude:sanity}
在相信任何计算结果之前：
\begin{enumerate}
  \item 将每个输入量取到其\omterm{def:g10:orders-of-magnitude:oom}{数量级}；
  \item 仅用 10 的幂重新计算，利用
        \cref{prop:g10:orders-of-magnitude:powers}，此步只需
        几秒钟；
  \item 比较：精确结果与估算结果应在
        十倍左右的范围内相符，且其\omterm{def:g10:orders-of-magnitude:unit}{单位}必须正确。若
        任一项检验未通过，须先查明错误，
        然后再继续。
\end{enumerate}
\end{method}

\begin{remark}[何必费心]\label{rem:g10:orders-of-magnitude:sanity}
计算器会欣然报告，一个苹果的质量为
$\qty{4e8}{kg}$——它无从知道更好的答案。没有任何公式能标出
荒谬之处，只有习惯能做到。物理学家先估算、后计算：
估算能发现小数点错位、\omterm{def:g10:orders-of-magnitude:unit}{单位}互换和
分数颠倒这三类错误，它们是各个层次上
多数错误答案背后的原因。
\end{remark}

\section{测量与不确定度}

没有任何仪器能读出真值。测量只有在
同时给出其可能偏离的程度时，才是可靠的。

\begin{definition}[绝对与相对不确定度]\label{def:g10:orders-of-magnitude:uncertainty}
设\omterm{def:g10:orders-of-magnitude:unit}{物理量} $x$ 的测量结果表示为
\[
x = x_{\text{m}} \pm \Delta x ,
\]
其中 $x_{\text{m}}$ 为测量值，
\emph{绝对不确定度}\index{不确定度!绝对} $\Delta x > 0$，与测量值\omterm{def:g10:orders-of-magnitude:unit}{单位}相同，界定可能的误差：
真值被认为位于
$x_{\text{m}} - \Delta x$ 与 $x_{\text{m}} + \Delta x$ 之间。
\emph{相对不确定度}\index{不确定度!相对}即比值
$\Delta x / x_{\text{m}}$，通常以百分数给出；它正是
“精确”一词所衡量的对象。
\end{definition}

\begin{method}[刻度仪器的读数]\label{met:g10:orders-of-magnitude:reading}
适用于直尺、量筒、表盘等：
\begin{enumerate}
  \item 读取最靠近指示位置的刻度，视线与
        刻度面垂直（斜视会使读数偏移）；
  \item 取不小于最小分度值一半的数作为\omterm{def:g10:orders-of-magnitude:uncertainty}{绝对
        不确定度}——若物体边缘或液面
        模糊，则应取更大值；
  \item 将结果写成 $x_{\text{m}} \pm \Delta x$ 的形式，并附\omterm{def:g10:orders-of-magnitude:unit}{单位}，
        且不保留超出不确定位的数字。
\end{enumerate}
\end{method}

\begin{omfigure}
\begin{tikzpicture}[x=0.9cm, font=\small]
  \fill[omThm, opacity=0.15] (6.5,-0.85) rectangle (7.5,1.45);
  \draw (-1.2,0) rectangle (10.6,1.1);
  \foreach \x in {0,...,10} \draw (\x,1.1) -- (\x,0.72);
  \foreach \x/\l in {0/120, 5/125, 10/130}
    \node[font=\footnotesize] at (\x,0.4) {\l};
  \fill[gray!25] (-1.0,-0.75) rectangle (7.0,-0.12);
  \draw (-1.0,-0.75) rectangle (7.0,-0.12);
  \draw[dashed, omThm] (7,-0.85) -- (7,1.45);
  \node[anchor=west, omThm] at (7.7,-0.5)
    {$L = (127.0 \pm 0.5)\,\unit{mm}$};
\end{tikzpicture}

{\small 读取毫米刻度直尺：杆的末端落在最靠近
$\qty{127}{mm}$ 的刻度处，阴影带覆盖两侧各半个
分度，它即为可信的结论。}
\end{omfigure}

\begin{example}[相对精度]\label{ex:g10:orders-of-magnitude:reading}
上述读数的\omterm{def:g10:orders-of-magnitude:uncertainty}{相对不确定度}为
$\frac{0.5}{127} \approx 0.4\%$。用同一把直尺测量一个
$\qty{12}{mm}$ 的螺栓，得到 $\frac{0.5}{12} \approx 4\%$：同一
仪器用于小物体时精度低了十倍。精度是
\emph{测量}的属性，而不是工具的属性。
\end{example}

\begin{remark}[绝对不确定度小而相对不确定度大——反之亦然]\label{rem:g10:orders-of-magnitude:precision}
激光测距站测得地月距离 $\qty{3.84e8}{m}$，精确到
约 $\pm\qty{3}{cm}$：\omterm{def:g10:orders-of-magnitude:uncertainty}{相对不确定度}为 $10^{-10}$，属
有史以来最精确的测量之一——其绝对误差
小到连尺子也不会夸耀。应始终追问：\emph{相对于什么}？数对
$(\Delta x,\ \Delta x / x_{\text{m}})$ 讲述了完整的情况，而
单独任何一个数都可能产生误导。
\end{remark}

\section{练习}

\begin{exercise}[$\star$]\label{exo:g10:orders-of-magnitude:1}
将下列数值换算为米，并用\omterm{def:g10:orders-of-magnitude:scinot}{科学记数法}表示：
\begin{enumerate}
  \item $\qty{750}{nm}$（红光的波长）；
  \item $\qty{2.5}{km}$;
  \item $\qty{3.2}{\micro m}$（一种土壤细菌）；
  \item $\qty{0.15}{Tm}$（地球与太阳的距离，
        以一种不寻常的装扮出现）。
\end{enumerate}
\end{exercise}

\begin{exercise}[$\star$]\label{exo:g10:orders-of-magnitude:2}
用\omterm{def:g10:orders-of-magnitude:scinot}{科学记数法}写出下列数值，并给出\omterm{def:g10:orders-of-magnitude:sigfig}{有效数字}的
位数：
\begin{enumerate}
  \item 地球与太阳的距离 $\qty{149600000}{km}$，以米为\omterm{def:g10:orders-of-magnitude:unit}{单位}；
  \item 氢原子的直径 $\qty{0.000000000106}{m}$；
  \item 世界人口约为 $\num{8100000000}$；
  \item 时长为 $\qty{0.00520}{s}$。
\end{enumerate}
\end{exercise}

\begin{exercise}[$\star$]\label{exo:g10:orders-of-magnitude:3}
不用计算器计算，并将每个结果
用\omterm{def:g10:orders-of-magnitude:scinot}{科学记数法}表示：
\[
(3\times10^{8})\times(2\times10^{-5}), \qquad
\frac{8\times10^{4}}{4\times10^{-3}}, \qquad
(5\times10^{6})^{2}, \qquad
\frac{(4\times10^{-3})\times(3\times10^{-9})}{6\times10^{-7}} .
\]
\end{exercise}

\begin{exercise}[$\star$]\label{exo:g10:orders-of-magnitude:4}
给出下列各量的\omterm{def:g10:orders-of-magnitude:oom}{数量级}，并附上\omterm{def:g10:orders-of-magnitude:unit}{单位}：人的身高
$\qty{1.7}{m}$；珠穆朗玛峰 $\qty{8848}{m}$；大肠杆菌（\emph{E.~coli}）
$\qty{2e-6}{m}$；地球的质量
$\qty{5.97e24}{kg}$。
\end{exercise}

\begin{exercise}[$\star$]\label{exo:g10:orders-of-magnitude:5}
一把尺子的刻度为毫米。用它测量一支铅笔，
得到 $\qty{127}{mm}$。
\begin{enumerate}
  \item 给出测量结果及其\omterm{def:g10:orders-of-magnitude:uncertainty}{绝对不确定度}。
  \item 计算\omterm{def:g10:orders-of-magnitude:uncertainty}{相对不确定度}。
  \item 用同一把尺子测量一块 $\qty{12}{mm}$ 的橡皮，
        两次测量中哪一次更精确？精确多少倍？
\end{enumerate}
\end{exercise}

\begin{exercise}[$\star\star$]\label{exo:g10:orders-of-magnitude:6}
利用 $c = \qty{3.00e8}{m/s}$ 与 $\qty{1}{ly} = \qty{9.47e15}{m}$：
\begin{enumerate}
  \item 将比邻星和天狼星的距离用米表示：
        比邻星（$\qty{4.24}{ly}$）、天狼星（$\qty{8.6}{ly}$）。
  \item \emph{旅行者一号}探测器距地球约 $\qty{2.4e13}{m}$，
        其无线电信号以速率 $c$ 传播，需要多长时间才能到达我们？
        结果用小时表示。
\end{enumerate}
\end{exercise}

\begin{exercise}[$\star\star$]\label{exo:g10:orders-of-magnitude:7}
求光从月球到达地球所需的时间，
月球距离 $\qty{3.84e8}{m}$，太阳距离 $\qty{1.496e11}{m}$，以及
海王星距离 $\qty{4.50e12}{m}$。每个答案用合适的\omterm{def:g10:orders-of-magnitude:unit}{单位}给出，并
补全句子：“当你看到海王星时，你看到的是它
当时的样子……”
\end{exercise}

\begin{exercise}[$\star\star$]\label{exo:g10:orders-of-magnitude:8}
用毫米刻度尺测得一令（$500$ 张）纸的厚度为
$\qty{52.0}{mm}$，不确定度为 $\pm\qty{0.5}{mm}$。
\begin{enumerate}
  \item 推算一张纸的厚度，
        并给出其\omterm{def:g10:orders-of-magnitude:uncertainty}{绝对不确定度}。
  \item 该厚度的\omterm{def:g10:orders-of-magnitude:uncertainty}{相对不确定度}是多少？
  \item 为什么用同一把尺子测量整令纸，比直接测量一张纸
        要好得多？请定量说明。
\end{enumerate}
\end{exercise}

\begin{exercise}[$\star\star$]\label{exo:g10:orders-of-magnitude:9}
取 $\qty{2e-10}{m}$ 为一个原子的直径。将原子一个挨一个地排列，
需要多少个原子才能跨越：一个长 $\qty{2.0}{\micro m}$ 的细菌；一根
长 $\qty{70}{\micro m}$ 的头发；一个 $\qty{1.7}{m}$ 高的人？
最后一个答案也请用\omterm{def:g10:orders-of-magnitude:oom}{数量级}表示。
\end{exercise}

\begin{exercise}[$\star\star$]\label{exo:g10:orders-of-magnitude:10}
某比例模型用一个足球代表太阳，太阳的直径为 $\qty{1.39e9}{m}$，
足球的直径为 $\qty{22}{cm}$。
\begin{enumerate}
  \item 计算该模型的比例因子。
  \item 求模型中地球的直径（实际直径 $\qty{1.27e7}{m}$），以及
        它与足球的模型距离（
        实际距离 $\qty{1.496e11}{m}$）。
  \item 模型中月球距模型地球多远（
        实际距离 $\qty{3.84e8}{m}$）？模型海王星
        （实际距离 $\qty{4.50e12}{m}$）又距离多远？
\end{enumerate}
\end{exercise}

\begin{exercise}[$\star\star$]\label{exo:g10:orders-of-magnitude:11}
\omterm{def:g10:orders-of-magnitude:sigfig}{有效数字}的应用。
\begin{enumerate}
  \item 一张 A4 纸的尺寸为 $\qty{21.0}{cm}$ 乘 $\qty{29.7}{cm}$，
        求其面积，并保留正确的\omterm{def:g10:orders-of-magnitude:sigfig}{有效数字}位数。
  \item 一名短跑运动员跑完 $\qty{100.0}{m}$ 所用时间为
        $\qty{12.3}{s}$，求平均速率，并正确取位。
  \item 用一句话说明，将面积报告为
        $\qty{623.70}{cm^2}$ 有何不诚实之处。
\end{enumerate}
\end{exercise}

\begin{exercise}[$\star\star\star$]\label{exo:g10:orders-of-magnitude:12}
一滴体积为 $\qty{1.0}{mm^3}$ 的油滴落在平静的水面上，
并扩散成直径为 $\qty{1.1}{m}$ 的圆形油膜——后面各章
将说明这层油膜恰好只有一个分子厚。
\begin{enumerate}
  \item 求油膜的面积，再求其厚度。
  \item 推算一个分子大小的\omterm{def:g10:orders-of-magnitude:oom}{数量级}，并
        将结果与阶梯上原子这一级进行核对。
\end{enumerate}
（这一粗略估算式的实验，最早由本杰明
·富兰克林在伦敦的一座池塘上尝试，是人类最早对分子世界作出的可靠测量之一。
）
\end{exercise}

\begin{exercise}[$\star\star\star$]\label{exo:g10:orders-of-magnitude:13}
估算一个人一生中的心跳次数。写出你的假设（静息心率、
寿命），运算时只用十的幂次和一位\omterm{def:g10:orders-of-magnitude:scinot}{尾数}，
并且仅给出最终答案的\omterm{def:g10:orders-of-magnitude:oom}{数量级}。为何
将其写成四位\omterm{def:g10:orders-of-magnitude:sigfig}{有效数字}
毫无意义？
\end{exercise}

\begin{exercise}[$\star\star\star$]\label{exo:g10:orders-of-magnitude:14}
一张纸厚 $\qty{0.10}{mm}$，每折叠一次，纸堆的厚度
便加倍。取 $2^{10} \approx 10^{3}$：
\begin{enumerate}
  \item 纸堆折叠多少次后，高度将超过珠穆朗玛峰
        （$\qty{8848}{m}$）？
  \item 折叠多少次后，纸堆将到达月球
        （$\qty{3.84e8}{m}$）？
  \item 实际上，一张纸至多只能折叠约七次。
        这是否使运算本身失效，
        还是只说明实验有问题？
\end{enumerate}
\end{exercise}

\begin{exercise}[$\star\star\star$]\label{exo:g10:orders-of-magnitude:15}
激光测距给出地月距离为
$\qty{3.84e8}{m} \pm \qty{3}{cm}$；一把好尺子测得桌子长为
$\qty{1.000}{m} \pm \qty{1}{mm}$。
\begin{enumerate}
  \item 分别计算两者的\omterm{def:g10:orders-of-magnitude:uncertainty}{相对不确定度}。哪一个测量
        更精确？二者相差多少倍？
  \item 为使对桌子的测量达到与月球测量相同的相对精度，
        \omterm{def:g10:orders-of-magnitude:uncertainty}{绝对不确定度}应为多大？
        将该长度与原子的大小进行比较，并
        恰如其分地作出结论。
\end{enumerate}
\end{exercise}

\section{问题：从原子到仙女座星系}

\begin{problem}[{周末问题——一粒沙子同时确定原子、太阳系、
银河系与可观测宇宙的尺度，并
确定你在这些尺度之间的阶梯上的位置}]
\label{pb:g10:orders-of-magnitude:1}
一粒沙子是物理学中最微不足道的对象，而这个周末它
将加班工作：先数清它的原子，再把太阳缩小到
与其同样大小，再用步数丈量宇宙；然后让光来完成测量，
最后数一数\emph{你}身体中的原子。全题所用数据如下：
沙粒直径 $\qty{0.5}{mm}$；原子直径 $\qty{2.5e-10}{m}$；
沙子的密度 $\qty{2500}{kg/m^3}$；直径：太阳
$\qty{1.39e9}{m}$，地球 $\qty{1.27e7}{m}$，银河
$\qty{9.5e20}{m}$，可观测宇宙 $\qty{8.8e26}{m}$；与我们的距离：
月球 $\qty{3.84e8}{m}$，太阳 $\qty{1.496e11}{m}$，海王星
$\qty{4.5e12}{m}$，比邻星 $\qty{4.0e16}{m}$，仙女座
星系 $\qty{2.4e22}{m}$；$c = \qty{3.00e8}{m/s}$；一个水分子
质量为 $\qty{3.0e-26}{kg}$，包含三个原子。

\medskip
\noindent\textbf{第一部分——沙粒的结构。}
\begin{enumerate}
  \item 以米为\omterm{def:g10:orders-of-magnitude:unit}{单位}，用\omterm{def:g10:orders-of-magnitude:scinot}{科学记数法}写出沙粒直径和原子直径，
        并分别指出各自最接近哪个
        SI \omterm{def:g10:orders-of-magnitude:prefixes}{词头}。
  \item 沿沙粒的一条直径并排排列，可放多少个原子？
  \item 将沙粒建模为边长为 $\qty{0.5}{mm}$ 的立方体，内部堆满
        原子。估算沙粒中的原子数，以
        \omterm{def:g10:orders-of-magnitude:oom}{数量级}表示。
  \item 求该立方体的体积，再求沙粒的质量。
        以毫克为\omterm{def:g10:orders-of-magnitude:unit}{单位}表示质量。
  \item 检验模型是否合理：用沙粒的质量除以其原子数，
        并与已知的原子质量\omterm{def:g10:orders-of-magnitude:oom}{数量级}作比较（
        约几个 $10^{-26}\,\unit{kg}$）。结论如何？
\end{enumerate}

\medskip
\noindent\textbf{第二部分——把太阳看作一粒沙子。}
我们构造一个宇宙的比例模型，其中太阳就是那粒沙子：
模型中的每一尺寸等于实际尺寸乘以因子 $s$。
\begin{enumerate}[resume]
  \item 求比例因子 $s$。
  \item 在模型中，求地球的直径及其与
        模型太阳（即沙粒）的距离。第一部分所述世界中，哪一个物体
        的大小与模型地球大致相当？
  \item 模型中的海王星绕沙粒运行，轨道距离多远？哪件
        家具能容纳整个行星系统？
  \item 模型中的比邻星距离多远？
  \item 两粒沙子相距 $\qty{14}{km}$，其间几乎一无所有：
        用一两句话说明该模型揭示了银河系
        是何等空旷。
\end{enumerate}

\medskip
\noindent\textbf{第三部分——用光来测量。}
\begin{enumerate}[resume]
  \item 光穿过真实的沙粒需要多长时间？
  \item 光从月球到达我们需要多长时间，
        从太阳到达我们又需要多长时间？
  \item 由 $c$ 及一年的时长，以米为\omterm{def:g10:orders-of-magnitude:unit}{单位}重新计算\omterm{def:g10:orders-of-magnitude:lightyear}{光年}，
        并验证数据中比邻星的距离
        约相当于 $\qty{4.2}{ly}$。
  \item 将仙女座星系的距离换算为\omterm{def:g10:orders-of-magnitude:lightyear}{光年}。大约在何时
        进入你眼睛的这束光离开仙女座星系，那时
        地球上正在行走的是什么生物？
  \item 光在一纳秒内传播多远？一台时钟频率为
        $\qty{3}{GHz}$ 的处理器每三分之一纳秒开始一次运算，
        用一句话说明为何快速计算机
        必须是小型计算机。
\end{enumerate}

\medskip
\noindent\textbf{第四部分——你，位于原子与仙女座星系之间。}
\begin{enumerate}[resume]
  \item 人体大致由 $\qty{70}{kg}$ 的水组成。利用题中数据，
        估算水中一个原子的平均质量，再求
        人体内的原子个数。
  \item 将这一结果与沙粒中的原子数（第 3 问）
        以及可观测宇宙中约 $2\times10^{23}$ 颗恒星相比较。
        分别给出两个倍数。
  \item 在这一阶梯上，你离原子更近，还是
        离可观测宇宙更近？比较比值 $\text{（你）}/\text{（原子）}$ 与
        $\text{（宇宙）}/\text{（你）}$，取你的身高
        为 $\qty{1.7}{m}$。
  \item 求恰好位于阶梯中间的尺度 $x$，
        即位于原子核（$\qty{e-15}{m}$）与
        可观测宇宙之间，$x$ 须满足 $x / 10^{-15} = 10^{26.9} / x$，
        （取 $\qty{8.8e26}{m} \approx 10^{26.9}\,\unit{m}$）。太阳系中的哪
        一个天体恰好具有这一大小？
  \item 压轴问题。将可观测宇宙按第二部分的因子
        $s$ 缩小，使太阳成为一粒沙子。给出模型宇宙的
        直径，先以米、再以\omterm{def:g10:orders-of-magnitude:lightyear}{天文单位}给出，
        并用一句话作结：即使每颗恒星都变成一粒沙子，
        宇宙仍远超人类的一切尺度——唯有
        十的幂次才能描述它。
\end{enumerate}
\end{problem}
